\documentclass{fank1basedocument}

\usepackage{hyperref}
\usepackage{array} % to not justify text in table
\usepackage[normalem]{ulem} % for strikethrough text

\settitle{Formal Needs, Features, and Requirements}

\fancyhead[L]{ROCKO}
\fancyhead[R]{2024-2025}

\begin{document}

\tableofcontents

\addcontentsline{toc}{section}{Needs}
\section*{Needs}
Educator \#1
\begin{enumerate}
	\item The robots used in the course are not robust (loses 40\% every year)
	\item It is difficult and expensive to find replacement parts
	\item A scaled-down robot needs to be affordable (similar price range to previous models ($200-$400))
	\item Having too many sensors can make the robot too complicated
	\item Some sensor APIs aren’t approachable
	\item Students want to work on a system that looks advanced
	\item It takes too long (2 weeks) for students to learn the platform (includes robot, programming language, and APIs)
\end{enumerate}
Educator \#2
\begin{enumerate}
	\setcounter{enumi}{7}
	\item It is difficult to provide examples when students use different IDEs
	\item It is difficult to help students who use different methods for calling APIs
	\item The current robots are no longer being produced
	\item The motors are more precise than the encoders
	\item There is no built in support for wireless communication, leading to many solutions
	\item The cameras are unreliable
	\item Modular platforms make it easier to adjust the course
\end{enumerate}
Hobbyist \#1
\begin{enumerate}
	\setcounter{enumi}{14}
	\item In FIRST Robotics Competition, students can’t reimplement things to learn how they work
	\item Bad quality components are time consuming and difficult to debug
	\item Closed platform, restrictions on motors (in FIRST Robotics Competition)
	\item The step between basics and more complicated features is too big
	\item Most kits are meant for learning software, not electrical components
\end{enumerate}
Hobbyist \#2
\begin{enumerate}
	\setcounter{enumi}{19}
	\item Need to set up the entire control system to run quick tests
	\item Wants adding sensors to be straightforward
	\item Too easy to break robots if a sensor comes unplugged
	\item Had issues with wires coming unplugged on breadboard robot
\end{enumerate}
Student/Benefactor
\begin{enumerate}
	\setcounter{enumi}{23}
	\item Things didn’t always work in mobile robotics
	\item Difficult to set up
	\item Didn’t have the computing power for SLAM on the robot
	\item Robot price should be \$750-\$1000
	\item FIRST Robotics Competition robots are too heavy
	\item FIRST Robotics Competition robots are too large
\end{enumerate}

\addcontentsline{toc}{section}{Features}
\section*{Features}
\begin{center}
	\begin{tabular}{|p{.1\textwidth}|p{.7\textwidth}|p{.12\textwidth}|}
		\hline
		\textbf{Number} & \textbf{Feature} & \textbf{Needs}\\
		\hline
		1 & Ability to add motors and sensors & 14, 15, 17, 21\\
		\hline
		2 & Ability to remove motors and sensors & 3, 4, 14, 15, 17\\
		\hline
		3 & Ability to substitute motors and sensors (compatible specs) & 3, 10, 14, 15, 17\\
		\hline
		4 & Prototype robot price should be \$750-\$1000 & 27\\
		\hline
		5 & Scaled-down robot price should be \$200-\$400 & 3\\
		\hline
		6 & Software platform should have code to communicate with provided motors & 9\\
		\hline
		7 & Software platform should have code to communicate with provided sensors & 5, 9\\
		\hline
		8 & Software platform should handle demos & 7\\
		\hline
		9 & Software platform should include various levels of tutorials & 7, 18\\
		\hline
		10 & Users should be able to change the code implementation of their motors & 15\\
		\hline
		11 & Users should be able to change the code implementation of their sensors & 15\\
		\hline
		12 & APIs should be easy to learn & 5, 7\\
		\hline
		13 & Users should be able to use programming languages that are easy to learn & 5, 7\\
		\hline
		14 & Software platform should be unified & 8\\
		\hline
		15 & Robot should be robust & 1, 23, 24\\
		\hline
		16 & Ability to survive drop tests & 1, 23\\
		\hline
		17 & Ability to view sensor outputs when the robot is not enabled & 5, 7, 22\\
		\hline
		18 & Wires should not come unplugged easily & 23, 24\\
		\hline
		19 & Electrical component quality should be appropriate for hobbyists & 11, 13, 16, 24\\
		\hline
		20 & Robot should be straightforward to set up & 7, 25\\
		\hline
		21 & Robot should look advanced & 6\\
		\hline
		22 & Ability to implement advanced software features & 26\\
		\hline
		23 & One person should be able to lift the robot & 28\\
		\hline
		24 & Ability to communicate wirelessly & 12\\
		\hline
		25 & Ability to balance & 6\\
		\hline
		26 & Robot should work & 24\\
		\hline
		27 & Robot has an articulated leg & 6\\
		\hline
		28 & Robot has sensors to detect the environment & 6\\
		\hline
		29 & Robot should be smaller than a microwave when fully extended & 29\\
		\hline
	\end{tabular}
\end{center}

\addcontentsline{toc}{section}{Requirements}
\section*{Requirements}
\addcontentsline{toc}{subsection}{Hardware}
\subsection*{Hardware}
\begin{center}
	\begin{tabular}{|p{.1\textwidth}|p{.7\textwidth}|p{.12\textwidth}|}
		\hline
		\textbf{Number} & \textbf{Requirement} & \textbf{Features}\\
		\hline
		1 & Robot stays intact and functioning when it falls over or is dropped from 8 inches high & 15\\
		\hline
		2 & Robot is under 25 pounds & 23\\
		\hline
		3 & Robot has encoders or potentiometers on each joint for closed-loop control & 25\\
		\hline
		4 & Center of mass is directly above wheels for balancing & 25\\
		\hline
		5 & Robot has a motor driving a wheel on the left and right side & 25\\
		\hline
		6 & Robot’s core functionality (driving and balancing) is assembled and wired & 26\\
		\hline
		\sout{7} & \sout{Robot has a motor driving the shoulder joint on the left and right side} & \sout{27}\\
		\hline
		8 & Total cost of robot is \$750 & 4\\
		\hline
		9 & Robot will be no larger than 20 x 18 x 10 inches when extended & 29\\
		\hline
		10 & Robot has a battery that has enough current delivery capacity to run all motors and sensors simultaneously & 24\\
		\hline
		11 & Robot power supply can be disconnected when manually handled but will not disconnect during operation & 18\\
		\hline
		12 & Robot is capable of distributing power on 12V and 5V interfaces & 26\\
		\hline
		13 & Robot battery can supply idle operation for 2hr & 26\\
		\hline
	\end{tabular}
\end{center}
Requirement 7 was removed after testing with Alpha Bot because of time it would require for controls and to improve packaging for the tank mode.

\addcontentsline{toc}{subsection}{Software}
\subsection*{Software}
Since the end-user's primary interaction with this system is through modifying code in their own development environment, most of our requirements are nonfunctional.

\addcontentsline{toc}{subsubsection}{Functional Requirements}
\subsubsection*{Functional Requirements}

\paragraph{Use Cases}
\begin{center}
	\begin{tabular}{|p{.1\textwidth}|>{\raggedright\arraybackslash}p{.2\textwidth}|p{.5\textwidth}|p{.12\textwidth}|}
		\hline
		\textbf{Number} & \textbf{Name} & \textbf{Description} & \textbf{Features}\\
		\hline
		1 & Run demo & The user runs some provided demo code on the robot  & 8\\
		\hline
		2 & Do tutorial & The user completes a tutorial to learn the robot platform & 9\\
		\hline
		3 & View sensor values & The user views the current values of sensors for debugging/troubleshooting & 17\\
		\hline
		4 & Communicate with robot using joystick & The user wants to drive the robot & 24\\
		\hline
		5 & Communicate with robot using computer & The user wants to view various pieces of data from the robot on their computer, such as sensor data and visualizations & 24\\
		\hline
	\end{tabular}
\end{center}

\pagebreak

\addcontentsline{toc}{subsubsection}{Nonfunctional Requirements}
\subsubsection*{Nonfunctional Requirements}

\paragraph{Scenarios}
\begin{center}
	\begin{tabular}{|p{.1\textwidth}|>{\raggedright\arraybackslash}p{.7\textwidth}|p{.12\textwidth}|}
		\hline
		\textbf{Number} & \textbf{Scenario} & \textbf{Features}\\
		\hline
		1 & \large{\textbf{Substitute Motor/Sensor - Integrability}}\newline
				\textbf{Source of Stimulus:} User\newline
				\textbf{Stimulus:} Substitutes a motor/sensor of similar spec\newline
				\textbf{Artifact:} The robot code\newline
				\textbf{Environment:} Development time\newline
				\textbf{Response:} The rest of the software works with the changed component\newline
				\textbf{Response Measure:} The other nodes only need minor changes, like retuning PIDs or changing the node it’s subscribed to
			& 3, 28\\
		\hline
		2 & \large{\textbf{Add Motor/Sensor - Integrability}}\newline
				\textbf{Source of Stimulus:} User\newline
				\textbf{Stimulus:} Adds a motor/sensor\newline
				\textbf{Artifact:} The robot code\newline
				\textbf{Environment:} Development time\newline
				\textbf{Response:} The motor/sensor can be interacted with in the rest of the system\newline
				\textbf{Response Measure:} Only nodes that used the motor/sensor need to be modified
			& 2\\
		\hline
		3 & \large{\textbf{Remove Motor/Sensor - Integrability}}\newline
				\textbf{Source of Stimulus:} User\newline
				\textbf{Stimulus:} Removes a motor/sensor\newline
				\textbf{Artifact:} The robot code\newline
				\textbf{Environment:} Development time\newline
				\textbf{Response:} The motor/sensor can be removed from the code\newline
				\textbf{Response Measure:} The other nodes only need minor changes, like retuning PIDs or changing the node it’s subscribed to
			& 3, 28\\
		\hline
		4 & \large{\textbf{First Setup - Usability}}\newline
				\textbf{Source of Stimulus:} User\newline
				\textbf{Stimulus:} Sets up the robot for the first time\newline
				\textbf{Artifact:} The robot\newline
				\textbf{Environment:} Runtime\newline
				\textbf{Response:} Will be fully set up and ready to run sample code\newline
				\textbf{Response Measure:} Takes no more than two hours
			& 20\\
		\hline
		5 & \large{\textbf{Pre-Provided Code - Usability}}\newline
				\textbf{Source of Stimulus:} User\newline
				\textbf{Stimulus:} Opens development environment\newline
				\textbf{Artifact:} Platform\newline
				\textbf{Environment:} Development time\newline
				\textbf{Response:} There is code to run the provided motors/sensors\newline
				\textbf{Response Measure:} The user doesn’t need to write this code
			& 6, 7, 28\\
		\hline
	\end{tabular}
\end{center}

\pagebreak

\begin{center}
	\begin{tabular}{|p{.1\textwidth}|>{\raggedright\arraybackslash}p{.7\textwidth}|p{.12\textwidth}|}
		\hline
		6 & \large{\textbf{Modify motor/sensor implementation - Modifiability}}\newline
				\textbf{Source of Stimulus:} User\newline
				\textbf{Stimulus:} Wants to modify a motor/sensor implementation\newline
				\textbf{Artifact:} Robot code\newline
				\textbf{Environment:} Development time\newline
				\textbf{Response:} User can change implementation\newline
				\textbf{Response Measure:} The rest of the code won’t be affected
			& 10, 11\\
		\hline
		7 & \large{\textbf{Learning Platform - Usability}}\newline
				\textbf{Source of Stimulus:} User\newline
				\textbf{Stimulus:} Starts learning how to code the robot\newline
				\textbf{Artifact:} The platform\newline
				\textbf{Environment:} Development time\newline
				\textbf{Response:} User can learn the basics of the platform (any APIs, the programming language, etc)\newline
				\textbf{Response Measure:} Within 1 week
			& 12, 13\\
		\hline
		8 & \large{\textbf{Balancing - Performance}}\newline
				\textbf{Source of Stimulus:} Robot\newline
				\textbf{Stimulus:} Runs balancing code\newline
				\textbf{Artifact:} Control nodes\newline
				\textbf{Environment:} Runtime\newline
				\textbf{Response:} Sensor values are retrieved, and outputs are calculated and sent to motors\newline
				\textbf{Response Measure:} Fast enough for balancing to occur effectively (early guess is around 20ms)
			& 25\\
		\hline
		9 & \large{\textbf{Unified Platform - Usability}}\newline
				\textbf{Source of Stimulus:} User\newline
				\textbf{Stimulus:} Wants to develop code for the robot\newline
				\textbf{Artifact:} The platform\newline
				\textbf{Environment:} Development time\newline
				\textbf{Response:} Has an IDE and programming language in place\newline
				\textbf{Response Measure:} They can be used to perform all normal functions
			& 14\\
		\hline
	\end{tabular}
\end{center}

\pagebreak

\paragraph{Shall Statements}
\begin{center}
	\begin{tabular}{|p{.1\textwidth}|p{.7\textwidth}|p{.12\textwidth}|}
		\hline
		\textbf{Number} & \textbf{Shall Statement} & \textbf{Scenarios}\\
		\hline
		1 & When the user substitutes a motor/sensor of matching spec during development time, the code shall work with the changed component with only minor changes needed to other nodes & 1\\
		\hline
		2 & When the user adds a motor/sensor during development time, the user shall be able to use that motor/sensor in other parts of the code within an hour (this doesn’t include time to develop driver for motor/sensor) & 2\\
		\hline
		3 & When the user removes a motor/sensor during development time, the motor/sensor can be removed from the code with only nodes that used that motor/sensor requiring changes & 3\\
		\hline
		4 & When the user sets the robot up for the first time, the robot will be fully set up within two hours & 4\\
		\hline
		5 & When the user opens the development environment, the platform shall provide code to interact with the given motors and sensors. & 5\\
		\hline
		6 & When the user wants to modify the implementation of a motor/sensor, the rest of the code shall require little to no modifications to work with the new implementation & 6\\
		\hline
		7 & When the user starts learning to code the robot, the platform shall allow the user to learn the basics (any APIs, the programming language, the IDE) within a week & 7\\
		\hline
		8 & When the robot runs balancing code, the overall loop time (sensor value retrieval, calculations, and motor commands) shall take no longer than 20ms & 8\\
		\hline
		9 & When the user wants to develop code for the robot, the platform shall provide an IDE and programming language that can handle all robot functions & 9\\
		\hline
	\end{tabular}
\end{center}

\addcontentsline{toc}{subsection}{Miscellaneous}
\subsection*{Miscellaneous}
\begin{center}
	\begin{tabular}{|p{.1\textwidth}|p{.7\textwidth}|p{.12\textwidth}|}
		\hline
		\textbf{Number} & \textbf{Requirement} & \textbf{Features}\\
		\hline
		1 & Produced documentation complies with web accessibility standards (WCAG 2) & N/A\\
		\hline
		2 & Produced documentation is interpretable at a 9th grade reading level or below & N/A\\
		\hline
	\end{tabular}
\end{center}

\end{document}
